% FILENAME: paper/section-Particles.tex

\section{Particles: Topological Knots and the Mass Gap}
\label{sec:Particles}

Many attempts to create a unified field theory have failed because they cannot naturally generate matter. \citet{goenner2004history} notes that the classical attempts by Einstein, Weyl, and Eddington to build particles out of spacetime geometry collapsed, and that Einstein himself lamented ``the problem of matter and quanta makes the construction fall apart'' when describing his 5-dimensional metric. Modern, higher dimensional theories, such as Kaluza-Klein supergravity, fail to incorporate chiral fermions \citep{witten1981search}. Consequently, physicists usually bolt external matter fields onto the gravitational Lagrangian by hand.

With CGD, particles are not fundamental points, but are instead knots within the connection itself. The self-dual (SD) chiral \slTwoC sector gives rise to gravity, but topological defects within the anti-self-dual (ASD) chiral \slTwoC sector gives rise to particles, while the curvature within the ASD sector produces mass. We will show how the interaction between these two chiral halves within CGD produces inertia.

Conceptually, \citet{wheeler1955geons} was the first to explore the idea of particles arising from the geometry of the universe. \citeauthor{wheeler1955geons} argued, ``Space time is not an arena for physics, it is all of classical physics'' \citep{wheeler1957nature}. He tried to build the ``geon'' out of the metric tensor and the electrodynamic field, but this approach ultimately failed to produce fermions.

\citet{thooft1974magnetic} and \citet{belavin1975pseudoparticle} moved away from Abelian fields and worked with Yang-Mills theories. They found solutions that looked like particles. Then, \citet{belavin1975pseudoparticle} discovered the BPST instanton, which are lumps of pure Yang-Mills fields that exhibit topological stability, although the BPST instanton strictly require Euclidean space. In this paper we will show the existence of topologically stable particles within Lorentzian space.

One of the biggest concerns with producing particles from a \spinFourC ontology is that this seems to cause a conflict with the strong-force. That is, the Standard Model describes the strong force using \suThree, and \suThree cannot be embedded within \spinFourC. In this section, we will address this problem by showing that CGD mathematically produces three ``colors''. We also show that due to how CGD particles are defined as part of the topology and cannot be separated from the Urbantke metric, CGD possesses a geometric mass gap.

We will begin our investigation into particles by considering a topological defect $(F^{(R)}_{\mu\nu} \neq 0)$ and then ask what this defect looks like to a macroscopic observer. In the Einstein Field Equations $(G_{\mu\nu} = 8 \pi G T_{\mu\nu})$, the $T_{\mu\nu}$ term is the stress-energy tensor. We will assume, in order to reach a contradiction, that this topological defect possess no mass or energy, giving $T_{\mu\nu} = 0$. However, \citet{krasnov2011plebanski} showed that if $T_{\mu\nu} = 0$, then the ASD curvature must be completely flat. However, this is a contradiction, because we begin our proof by assuming that a knot exists by stating $F^{(R)}_{\mu\nu} \neq 0$. Therefore, we conclude
\begin{equation}
    F^{(R)}_{\mu\nu} \neq 0 \implies T_{\mu\nu} \neq 0 .
    \quad \leanref{CGD.Particles.dynamicMatterExistence}
    \label{eq:dynamicMatterExistence}
\end{equation}

In CGD, we must build particles out of \spinFourC somehow, but this leaves us with several options. \uOne is the simplest non-empty subgroup, but this immediately runs into problems because \uOne is commutative. When we plug a commutative field into the Urbantke Metric (Equation \ref{eq:urbantkeMetric}), we see that the volume collapses, which would mean that no particles experience time or even relative distance. This means
\begin{equation}
    [F_{\mu\nu}, F_{\rho\sigma}] = 0 \implies \det(g_{\mu\nu}) = 0 .
    \quad \leanref{CGD.Particles.kinematicSingleColorDegeneracy}
    \label{eq:kinematicSingleColorDegeneracy}
\end{equation}
\noindent More generally, if the field does not have enough independent Lie algebra generators, or ``colors'', to fill in all of the diagonals of the $3 \times 3$ matrix $g_{\mu\nu}$, then the volume will collapse. We can use this observation to show that particles require exactly three colors
\begin{equation}
    \det(g_{\mu\nu}) \neq 0 \implies \big( [F_{\mu\nu}, F_{\rho\sigma}] \neq 0 \big) \land \big( \forall a \in \{1,2,3\}, \, F^a_{\mu\nu} \not\equiv 0 \big) .
    \quad \leanref{CGD.Particles.kinematicThreeColorRequirement}
    \label{eq:kinematicThreeColorRequirement}
\end{equation}

Let us next consider defects in the ASD chiral sector. In topology, the Cartan-Maurer degree, or winding number, counts how many times a field wraps around a gauge group \citep{nakahara2003geometry}. Since \citet{krasnov2011plebanski} proved that only \suTwo and \slTwoC are capable of generating the Urbantke Metric, then for CGD we only need to only consider the ASD chiral \slTwoC gauge field, where the topological winding numbers are described by the \suTwo subgroup. If a particle exists in \suTwo, it is a non-Abelian knot, and \citet{belavin1975pseudoparticle} proved that knots in \suTwo cannot decay or unwrap into the trivial vacuum. We can extend these results to show that they also apply within CGD:
\begin{equation}
    \int_{S^3} \mathrm{Tr}\left(\mathcal{A} \wedge d\mathcal{A} + \frac{2}{3} \mathcal{A} \wedge \mathcal{A} \wedge \mathcal{A}\right) \neq 0 \implies \mathcal{A} \text{ cannot continuously decay to } 0.
    \quad \leanref{CGD.Particles.kinematicTopologicalStability}
    \label{eq:kinematicTopologicalStability}
\end{equation}

Next, we use the existence of the Lorentzian Witness (Equation \ref{eq:dynamicExactLorentzianSolution}) to perform an analysis that has not previously been possible. In \citet{belavin1975pseudoparticle}, particles strictly existed in an Euclidean field. However, since a Euclidean field does not possess a time dimension, these topological defects possess an `action' but they cannot possess a rest mass or energy. Since CGD dynamically generates a Lorentzian space, we evaluate a particle by evaluating the 3-dimensional Cauchy slice of spacetime at time $t$ where the spatial boundary $\partial \Sigma$ compactifies to $S^3$. We define the inertial mass as:
\begin{equation}
    m_{\text{inertial}} = \left| \int_{\partial \Sigma} \omega_{\text{CM}}(\mathcal{A}^{(R)}) \right| .
    \quad \leanref{CGD.Particles.inertialMass}
    \label{eq:inertialMass}
\end{equation}

The non-compact boost directions of the \slTwoC gauge field form a contractable space $(\mathbb{R}^3)$, so they possess a topological winding number of zero. This excludes negative-mass ghost states. Since the stable particle must possess a topological homeomorphism to the \suTwo boundary, the integrated charge must evaluate to a non-zero integer. This shows that the topological defect must possess a strictly positive mass gap:
\begin{equation}
    \mathcal{A}^{(R)}\big|_{\partial \Sigma} \cong \text{SU}(2) \implies m_{\text{inertial}} > 0 .
    \quad \leanref{CGD.Particles.topologicalMassGap}
    \label{eq:topologicalMassGap}
\end{equation}

It's important to emphasize that this is not the Millennium Prize Problem \citep{jaffe2006quantum}. The proof for Equation \ref{eq:topologicalMassGap} relies on the global topological mapping degree of a continuous classical boundary, and it's only possible because the Urbantke metric is inseparable from the topology. While we are able to show that a mass gap exists in CGD, this proof shows nothing about the similar problem of proving a mass gap for any theory, such as the Standard Model, that possesses an independent metric background.

Now that we have characterized the mass for \suTwo knots, we can also look at simpler Abelian \uOne projections of the ASD connection. To do this, we will use the Abelian field strength tensor which incorporates both the standard curl and the non-Abelian commutator:
\begin{equation}
    F_{\mu\nu}^{(U(1))} = \partial_\mu A_\nu - \partial_\nu A_\mu + [A_\mu, A_\nu] \quad \leanref{CGD.Particles.abelianFieldStrength}
\end{equation}
\noindent By applying the Duan-Ge topological decomposition \cite{duan1979su2}, we can extract the \uOne electric current by taking the dual divergence of the field strength:
\begin{equation}
    J^\mu = \epsilon^{\mu\nu\rho\sigma} \partial_\nu F_{\rho\sigma} \quad \leanref{CGD.Particles.emergentElectricCurrent}
\end{equation}
\noindent Since the \spinFourC universe is a smooth manifold, the $n$-dimensional Clairaut's theorem forces the partial derivatives to commute \citep{rudin1976principles}. This forces the divergence of the current to vanish, proving charge conservation:
\begin{equation}
    \partial_\mu J^\mu = 0
    \quad \leanref{CGD.Particles.topologicalChargeConservation}
\end{equation}

A classical spinor field traditionally commutes, meaning it obeys bosonic statistics. As noted earlier, this has historically been a problem for classical theories. However, while the local spinor in CGD is a commuting field, CGD mathematically possesses a macroscopic defect that behaves differently. In \citeyear{witten1983current}, \citet{witten1983current} demonstrated that topological solitons exhibit fermionic spin statistics. In \suThree geometries he used a Wess-Zumino term, but this term disappeared for \suTwo geometries, such as the one used in CGD. For the \suTwo case, we use the Finkelstein-Rubinstein theorem \citep{finkelstein1968connection}, which proves that an adiabatic $2\pi$ spatial rotation of a degree-1 \suTwo soliton traces a non-trivial loop in configuration space $(\pi_4(\text{SU}(2)) = \mathbb{Z}_2)$. This proves that macroscopic topological defects in CGD possess fermionic spin statistics:
\begin{equation}
    \text{Weight}(U_{\text{rot}}) = -1 .
    \quad \leanref{CGD.Particles.kinematicMacroscopicFermion}
    \label{eq:kinematicMacroscopicFermion}
\end{equation}

